% V. Aumento de temperaturas globais; V. Ação humana no clima; V. Sul Global mais afetado proporcionalmente
% 3. Aumento de heat waves -> Dunn;
% 4. Aumento de precipitação;
% 5. Aumento de Enchentes (precipitação/precipitação acumulada);
% 6. El Nino e La Nina

% Chapter 2.3.1.1.3 (Temperatures during the instrumental period - surface) of the AR6 WGI;
As últimas décadas foram marcadas por quebras constantes de recordes climáticos, especialmente em máximas de
temperatura, evidenciando de forma empírica o impacto das mudanças climáticas em escala global. As anomalias de
temperatura, definidas como a diferença entre a temperatura observada em um determinado período e uma temperatura média
de referência calculada para um período base, são amplamente utilizadas como indicadores de tendências e mudanças de
longo prazo nos padrões climáticos, pois podem capturar variações consistentes e sistemáticas não aparentes em dados
absolutos, permitindo uma análise mais clara das mudanças globais. Nas últimas décadas, as anomalias de temperatura
apresentam um padrão de crescimento constante, como pode evidenciado na Figura~\ref{temp_anomaly}.

\begin{figure}[htb]
	\caption{\label{temp_anomaly} Anomalia de Temperatura referente ao período de 1850 a 2024}
	\begin{center}
	    \includegraphics[scale=0.455]{img/Temperature Anomaly.png}
	\end{center}
	\legend{Fonte: \textit{National Oceanic and Atmospheric Administration} (NOAA) \cite{ncei_noaa}}
\end{figure}

De acordo com o sexto - e mais recente - relatório do Painel Intergovernamental sobre Mudanças Climáticas (IPCC),
estabelecido pelas Nações Unidas e amplamente reconhecido como a principal referência na análise do conhecimento
científico relacionado às mudanças climáticas, a superfície terrestre tem experiênciado um aumento expressivo de
temperatura em comparação com períodos históricos. Considerando o intervalo de 1850 a 1900 (utilizado como base
pré-industrial por ser o primeiro período com boas observações de temperaturas terrestres registradas), a temperatura
média da superfície terrestre\footnote{Dados referentes à medida do GMST - Temperatura Média Global da Superfície - que
representa uma média ponderada entre as temperaturas das massas terrestres e das massas oceânicas.} no período de 2001 a
2020 foi aproximadamente 1.07ºC mais elevada. Além disso, a tendência de aumento da temperatura somente nas últimas
quatro décadas (1980 a 2020) é de 0,76ºC, evidenciando uma aceleração no aquecimento global \cite{gulev}.

Esse aumento, aparentemente pequeno, não reflete a real magnitude das modificações no sistema terrestre como um todo. É
importante notar que um aumento de 1ºC representa uma média global; portanto, houveram locais no qual o aquecimento foi
significativamente mais pronunciado, como nos polos. Nessas regiões, mudanças pequenas na temperatura podem levar ao
derretimento de geleiras e, logo, na diminuição do albedo da Terra, isto é, a capacidade de refletir os raios solares de
volta para o espaço. Portanto, nessa situação, a Terra absorverá mais calor, levando à ainda mais derretimento das
geleiras, criando assim um ciclo de retroalimentação. Além disso, é relevante destacar que um aquecimento de 1ºC em toda
a superfície terrestre é extremamente significativo devido à massa do planeta. Para que um aumento dessa magnitude
ocorra em um corpo com a massa da Terra, a quantidade de energia térmica\footnote{De forma bastante simplificada, essa
energia pode ser vista pela fórmula $ Q = mc\Delta T $, onde $m$ é a massa, $c$ é o calor específico, e $\Delta T$ a
variação de temperatura.} necessária é imensa.

% Chapter 3.3 (Human Influence on the Atmosphere and Surface) of the AR6 WGI;
É indubitável que a ação humana foi o maior contribuinte para o Aquecimento Global \cite{ipcc, gillett}. As atividades
industriais que sustentam os padrões de consumo modernos são as principais emissoras de gases responsáveis pelo efeito
estufa, como dióxido de carbono (\ce{CO2}), metano (\ce{CH4}), óxido nitroso (\ce{N2O}), entre outros. Esses gases
possuem a capacidade de absorver a radiação infravermelha emitida pela superfície da Terra, retendo o calor que, de
outra forma, seria dissipado no espaço. 

Embora o efeito estufa seja considerado um fenômeno essencial graças à concentração natural de certos GEEs (sem ele a
temperatura média da superfície terrestre seria significativamente menor), a ação humana tem intensificado esse processo
alarmantemente. Desde o início da Revolução Industrial, a concentração de GEEs na atmosfera aumentou gradualmente,
principalmente devido à queima de combustíveis fósseis. Estima-se que, entre 1850 e 2019, as emissões acumuladas de
\ce{CO2} foram de 2400±240 Gt\ce{CO2} \cite{friedlingstein}, resultando em uma concentração de \ce{CO2} comparavél com o
início do máximo térmico do Paleoceno-Eoceno \cite{gingerich}, um período associado a uma das mais recentes extinções em
massa.

Esses dados destacam o papel intensificado das atividades humanas no agravamento do efeito estufa e evidenciam a
necessidade urgente de iniciativas coordenadas em escala global para mitigar os danos climáticos e garantir um futuro
mais sustentável. Ações como o Acordo de Paris são essenciais para garantir a coordenação das ações climáticas ao
estabelecer metas e responsabilidades de cada nação. Firmado em 2015 pelas nações participantes do Convenção-Quadro das
Nações Unidas sobre a Mudança do Clima (UNFCCC), o acordo estabelece a necessidade de reduzir as emissões de GEEs em
escala global, visando manter o aumento de temperatura abaixo de 2ºC em relação à base pré-industrial, além de
providenciar mecanismos para que nações em desenvolvimento possam mitigar os efeitos das mudanças climáticas e se
adaptar a elas sem comprometer severamente suas atividades econômicas essenciais.

% Chapter 8.2 of the AR6 WGII and Chapter 12.3.6 of the AR6 WGII;
\subsection{Contribuições e vulnerabilidades climáticas}
\label{vul}
A América Latina é um grande contribuinte na emissão de GEEs, contabilizando aproximadamente 11\% das emissões
históricas globais de \ce{CO2} \cite{friedlingstein}. Contudo, a maior parte dessa emissão vem na forma de uso de terra,
como agricultura, pecuária e desmatamento. Esse padrão é característico de muitas regiões em desenvolvimento, onde a
produção industrial está voltada predominantemente para atender às demandas de exportação de bens e matéria-prima, em
vez de ser direcionada ao consumo interno. Consequentemente, as contribuições dessas regiões para o aquecimento global
vem com o papel de fornecedores globais, beneficiando outras partes do mundo, enquanto enfrentam os desafios associados
a essas emissões. 

Além disso, essas regiões são consideravelmente mais vulneráveis às consequências das mudanças climáticas (REFERÊNCIA),
não apenas por se localizarem em áreas tipicamente mais tropicais, onde o aumento das temperaturas é mais intenso, mas
também devido à infraestrutura limitada para prevenir e mitigar desastres. Uma ampla gama de fatores determina a
resiliência de uma região em relação as mudanças climáticas, como desenvolvimento desigual, fracas instituições e
pobreza. Comunidades mais vulneráveis são afetadas desproporcionalmente por efeitos em, por exemplo, perdas no setor
básico da agricultura.

% -> Vulnerabilidade dos que contribuem menos. Tentar ao máximo ficar no tópico de América Latina;

% IPCC AR6 WGI Chapter 11.4 (Heavy Precipitation) and Chapter 11.5 (Floods); TODO IMAGEM!
\subsection{Mudanças nos padrões de precipitação}
\label{rain}
Mudanças em média de temperatura são processos graduais que se tornam difícies de perceber em espaços curtos de tempo. A
comparação de temperatura entre um ano e o anterior não mostrará a tendência real das mudanças, especialmente
considerando efeitos naturais de variabilidade climática, como o El Ñino e a La Ñina. No entanto, essas mudanças
graduais afetam diretamente nos eventos climáticos extremos facilmente perceptíveis. Esse aumento da temperatura, coloca
em desequilibrio delicados mecanismos atmosfericos levando à mudanças em, por exemplo, aumento de temperaturas máximas,
redução de temperaturas mínimas e ondas de calor\footnote{Definido por Dunn et al. (2024) \cite{dunn} como 3 ou mais
dias no qual a temperatura se encontra acima do 90º percentil}.

Um evento climático cujo mudança se destaca é a precipitação. Em essência, o aumento de temperatura leva à uma maior
evaporação da água e também a maior capacidade atmosférica de reter esse vapor de água, o que resulta em chuvas mais
intensas e irregulares. Uma definição mais completa pode ser dada por Seneviratne et al. (2021) no capítulo 11.1 do
sexto relatório, segundo grupo, do IPCC (AR6 WGII): 

\begin{citacao}
As mudanças na temperatura também controlam as variações no vapor d'água por meio do aumento da evaporação e da
capacidade de retenção de água da atmosfera. Em escala global, o conteúdo de vapor d'água integrado na coluna
atmosférica aumenta aproximadamente de acordo com a relação de Clausius–Clapeyron, com um aumento de cerca de 7\% para
cada 1°C de aquecimento médio global da superfície. (...) Projeções de modelos climáticos mostram que o aumento do vapor
d'água resulta em elevações significativas nos extremos de precipitação em todas as regiões, com uma magnitude que varia
entre 4\% e 8\% por 1°C de aquecimento da superfície.

\textit{Seneviratne et al. (2021)\cite{ipcc_ar6_wg2_chap11}, traduzido pelo autor.}

\end{citacao}

% definir alta precipitação;
É importante notar que o aumento das temperaturas não está somente associado a um incremento na precipitação total, mas
também à intensificação dos extremos de precipitação. Observa-se não só um aumento no volume máximo de chuva em um único
dia, evidenciando a maior intensidade dos eventos, como também no acumulado de precipitação em períodos de cinco dias,
indicando uma maior duração desses extremos \cite{dunn}.

Contudo, mesmo que a tendência global aponte para um aumento de precipitação extrema, isso não necessariamente implica
que todas as regiões do planeta sentirão este efeito de forma uniforme, pois acompanhado de um aumento nas chuvas, está
uma mudança nos padrões da mesma. Os complexos mecanismos atmosféricos que regulam a precipitação são influenciados
pelas mudanças graduais de temperatura, podendo tanto aumentar quanto reduzir não apenas a quantidade, mas também a
intensidade das chuvas. Essas mudanças na precipitação estão diretamente relacionadas à um aumento na quantidade de
secas e de enchentes, pois, se a distribuição de precipitações aumenta ou diminui, os limiares que levam à esses eventos
são mais facilmente alcançados. Essas variações são específicas a cada região, dependendo das particularidades
climáticas locais, tornando difícil a definição de uma tendência que será sentida de forma uniforme.

No Brasil, por exemplo, esses dois efeitos são observados de forma simultânea. No Nordeste, a região semiárida mais
populosa do mundo, há uma redução na frequência e intensidade das chuvas, o que tem levado a um aumento na ocorrência de
secas mais severas e frequentes, um problema crônico na área. Além disso, projeções indicam que o Nordeste enfrentará
aumentos de temperatura mais intensos do que outras regiões \cite{hoegh-guldberg}, agravando ainda mais os impactos da
seca (\autoref{vul}). Por outro lado, na região Sudeste, observa-se um aumento nas chuvas intensas desde o início do
século XX \cite{dunn}, e, acompanhado das mesmas, uma intensificação em enchentes e deslizamentos, que são a maior causa
de morte relacionada à desastres naturais na região \cite{ufsc}. 

% IPCC AR6 WGI - Chapter 8
\subsection{Enchentes e suas tendências}
Enchentes são caracterizadas pelo aumento do nível da água em áreas que normalmente permanecem secas. Esse acúmulo
prolongado ocorre devido à dificuldade de escoamento, seja pela incapacidade do solo de absorvê-la ou pela lentidão na
drenagem natural, como o fluxo em grandes corpos d'água. Em geral, elas podem ser classificadas em três tipos
principais: enchentes pluviais, resultantes de chuvas intensas que excedem a capacidade de absorção do solo, causando
alagamentos; enchentes fluviais, provocadas pelo transbordamento de corpos d'água – geralmente rios – quando o nível da
água ultrapassa suas margens, frequentemente devido a chuvas persistentes, derretimento de gelo ou acúmulo de
sedimentos; e enchentes costeiras, que ocorrem em áreas litorâneas devido ao aumento do nível do mar, grandes ondas ou
fenômenos como ciclones e tsunamis.

As causas de uma enchente nem sempre podem ser facilmente atribuídas a um único evento específico, pois geralmente
dependem de uma combinação de fatores. No caso das enchentes fluviais, por exemplo, não é apenas a ocorrência de chuvas
intensas no corpo d'água principal ou em seus tributários que provoca o transbordamento. Se o corpo hídrico tiver
capacidade de escoar a água de forma eficiente, o alagamento pode ser evitado. Muitos corpos d'água, cuja principal
fonte de abastecimento são enchentes pluviais, possuem uma morfologia adaptada para lidar com essa condição. Outros
fatores que contribuem para o transbordamento incluem a área de escoamento disponível em canais naturais e artificiais,
a umidade prévia do solo, sua capacidade de percolação, a presença de ventos fortes e sua interação com o fluxo da água,
a taxa de evaporação após o evento, a capacidade de absorção da vegetação local, o declive do terreno em relação à área
de escoamento, entre outros. Assim, não se pode associar exclusivamente o aumento das chuvas intensas (\autoref{rain})
ao crescimento de todos os eventos extremos desse tipo (REFERÊNCIA).

% ######################################################################################################################
%EFEITOS GEOGRAFICOS QUE LEVAM A ENCHENTES -> MOFOLOGIA DE RIOS, ESCOAMENTO, INTERAÇÃO COM VENTOS, DECLIVE E CANAIS:

% ######################################################################################################################
Mesmo assim, as enchentes são um fenômeno natural e frequentemente característico de certas regiões. Áreas com invernos
rigorosos, marcados pela acumulação de neve, costumam enfrentar enchentes durante o verão, quando o rápido derretimento
do gelo transporta grandes volumes de água de uma região para outra, sustentando ecossistemas que dependem desse influxo
sazonal. Regiões influenciadas por monções — ventos sazonais que trazem para o interior um acúmulo de umidade e
precipitação marítima — também são marcadas por cheias em seus corpos d'água, essenciais para a agricultura local, como
ocorre nos tributários dos rios Ganges e Amarelo. De maneira semelhante, as cheias anuais regulares do Rio Nilo são
fundamentais para a fertilidade da região. 

Entretanto, mesmo sendo uma parte integral do ciclo da água, a presença de assentamentos humanos geralmente corresponde
com a existência de grandes corpos d'água que podem ser, dadas circunstâncias, vulneráveis à enchentes. Essas imundações
podem invadir esses locais de forma inesperada e veloz, colocando vidas em risco, comprometendo a segurança das
comunidades e causando danos estruturais que amplificam a disrupção das atividades humanas. Fatalidades causadas por
afogamentos são comuns em regiões alagadas, não apenas pela rapidez com que as enchentes ocorrem, mas também devido às
características do escoamento das águas, que podem gerar correntes fortes e imprevisíveis. As enchentes também criam
condições propícias para o agravamento de outras condições sanitárias. Elas podem transmitir doenças por meio da água
contaminada, além de comprometer sistemas críticos, como a contaminação de fontes de água potável, a interrupção do
fornecimento de energia elétrica, afetando hospitais e serviços essenciais, e a destruição de cadeias de suprimentos.
Esses fatores podem isolar comunidades, mesmo aquelas não diretamente afetadas pela inundação, dificultando o acesso a
mantimentos e recursos básicos.

Os efeitos das enchentes, no entanto, não se limitam às comunidades urbanas. Atividades econômicas fundamentais, como a
agricultura, estão entre as mais prejudicadas por esses desastres. A necessidade de grandes extensões de terra e a
proximidade com corpos d’água para irrigação tornam o setor agrícola particularmente vulnerável. Em situações de
enchente, a safra pode ser destruída não somente pela força da água, mas também pela permanência prolongada de água
nessas regiões, já que muitas espécies de plantas não conseguem sobreviver a tais condições. Além da destruição de
colheitas e da perda de gado, as enchentes frequentemente resultam em danos à infraestrutura essencial, como sistemas de
irrigação, equipamentos, instalações de processamento e transporte, além de alterações nas terras cultiváveis. De acordo
com a Organização de Alimentos e Agricultura das Nações Unidas \cite{onu}, enchentes são responsáveis por mais da metade
de todas as perdas em colheitas causadas por desastres naturais. Entre 2008 e 2018, essas perdas em países em
desenvolvimento ultrapassaram 21 bilhões de dólares, tornando as enchentes a segunda maior causa de danos econômicos no
setor agrícola, atrás apenas das secas. Em países em desenvolvimento, onde setores primários têm um peso significativo
na economia, os efeitos desses desastres são amplamente sentidos de forma indireta no enfraquecimento econômico, podendo
levar a uma maior situação de vulnerabilidade e despreparo.

\subsection{Vulnerabilidade urbana à enchentes}
\lipsum[1-5]

% Usos artificiais para enchentes -> Geração de energia, represas e defesa (Dutch Waterline, Stelling van Amsterdam);

% Falar aqui de trabalhos como o do Ta;

% El ÑINO:
% https://wmo.int/media/news/global-temperature-record-streak-continues-climate-change-makes-heatwaves-more-extreme